% =====================================================================
%  Chapter 28 — The \idxw{Portfolio} Decision
% =====================================================================

\chapter{The Portfolio Decision}
\label{ch:portfolio}

A platform is not justified by what it can do. It is justified by what it
does that nothing else does, at a cost the decision can bear, against a
requirement that is written down. This chapter is the book doing that
arithmetic, and it is written in the form a programme would use it.

\section{The jobs to be done}

The accountability chain of Chapter 10 defines the work. Each link needs a
measurement, and each measurement needs an instrument. Naming the jobs
first is what stops the comparison from becoming an argument about
instruments.

\begin{table}[!htbp]
\centering
\small
\caption{The jobs, and what each one is worth.}
\label{tab:jobs}
\begin{tabularx}{\textwidth}{@{}P{0.78in}P{1.78in}Y@{}}
\toprule
\rowcolor{panel}
\textbf{Job} & \textbf{Deliverable} & \textbf{What it is worth} \\
\midrule
J1 & a diameter, without a thermal model & removes a factor of 1.68 in
energy \\
\rowcolor{panel}
J2 & a density & removes a factor of 2.69 in energy \\
J3 & a resolved shape and a pole & the volume, and the spin state that
YORP depends on \\
\rowcolor{panel}
J4 & a mass & the density, and with it J2 \\
J5 & an energy & the hazard class \\
J6 & an impact probability at a long horizon & the keyhole position \\
\bottomrule
\end{tabularx}
\end{table}

The jobs are not equal in value, and \tabref{tab:jobs} is how they are
ranked. J1 and J2 are worth a factor of 4.5 in energy between them. J6 is
worth a factor of 433 in the impact probability in the tail regime of
Chapter 14, and nothing at all in the near regime.

\section{The scorecard}

\begin{table}[!htbp]
\centering
\small
\caption{The portfolio scorecard. Each platform is scored against the jobs
it can close, the population it reaches, and the limit that binds it.}
\label{tab:scorecard}
\begin{tabularx}{\textwidth}{@{}P{1.11in}P{0.72in}P{0.88in}Y@{}}
\toprule
\rowcolor{panel}
\textbf{Platform} & \textbf{Jobs} & \textbf{Population} & \textbf{Binding
limit} \\
\midrule
Ground survey and lightcurve & J3 (partly) & thousands & the convex hull \\
\rowcolor{panel}
Ground adaptive optics & J1, J3 & a handful & close approach and phase
angle \\
Ground bistatic radar & J1, J2, J3, J4 & dozens & the fourth-power wall
and the radar-crossing selection \\
\rowcolor{panel}
NEO Surveyor and NEOMIR & J1 & the whole discovery stream & a model, and no
shape or mass \\
3.5 m space telescope & J3 (photometry) & survey plus rapid response & no
mid-infrared channel \\
\rowcolor{panel}
Low-Earth-orbit constellation & J1, J3 (partly) & proposed & orbit and
brightness \\
Lunar surface array & J1, J3, J4 at an encounter, J6 & 37 per year, on
request & the baseline, the dust and the night \\
\bottomrule
\end{tabularx}
\end{table}

The scorecard is not a ranking, because the jobs are not interchangeable.
It is a map of coverage, and it shows three things.

The thermal channel is covered, twice, by funded missions. Nothing needs to
be built there, and the proposal should not attempt it.

The shape channel is covered thinly and intermittently: by close-range
images, by radar for a small radar-crossing minority, and by convex
lightcurve inversions for everything else. This is the gap.

The mass channel is covered by radar, for a few dozen objects, and by
nothing else. It is the smallest gap in terms of population and the largest
in terms of consequence.

\section{The marginal-value argument}

The management question is not ``which platform is best'' but ``which
purchase removes the most error per unit of effort, and which purchase
removes an error that no other purchase removes''.

That framing changes the answer in three ways.

First, it eliminates the platforms that duplicate an existing capability.
The thermal channel is duplicated, so it costs nothing in marginal terms to
leave it where it is.

Second, it favours the platform that closes the dominant term. For the
energy chain, the dominant term in the small class is the density, and the
density is a mass measurement, which is a close-encounter measurement. A
platform that is not present at the encounter cannot deliver it at all.

Third, it makes the service model worth more than the optics. The alert
path from the risk list to the instrument costs almost nothing in
engineering terms, and the documented failure of 2024~YR4 was that it did
not exist.

\begin{table}[!htbp]
\centering
\small
\caption{The marginal-value comparison: what each purchase adds that nothing
else adds.}
\label{tab:marginal}
\begin{tabularx}{\textwidth}{@{}P{1.36in}P{1.20in}Y@{}}
\toprule
\rowcolor{panel}
\textbf{Purchase} & \textbf{Marginal contribution} & \textbf{Assessment} \\
\midrule
A thermal survey mission & none; two already exist & decline \\
\rowcolor{panel}
More ground adaptive optics & marginal; the constraint is opportunity, not
aperture & decline \\
A better radar & logarithmic in range & hold; it is indispensable for the
few dozen \\
\rowcolor{panel}
An alert path & the whole routing gap, at negligible cost & do it first, and
do it regardless of platform \\
A lunar shape array & the shape channel and the encounter & the proposal \\
\bottomrule
\end{tabularx}
\end{table}

\tabref{tab:marginal} is the honest form of the argument, and it contains
two things that a lunar advocate would rather not say. The first is that the
alert path should be built whether or not the lunar station is built, and
that it should be built first. The second is that a better radar is
indispensable and the lunar station does not replace it.

\begin{seamnote}{The reasoning that gets the Moon chosen, and the reasoning
that does not}
The Moon is chosen because of a conjunction of three facts that no other
platform shares. It is the only place where a kilometre-scale stable
baseline can be built, which is what J1 and J3 need. It is a platform that
can be present at a close encounter, which is what J2 and J4 need. And it is
the only vantage that reaches the Sun-ward discovery window, which is where
the objects still arriving at three thousand a year are found.

Any one of these alone would not justify it. The thermal channel alone would
argue against it. The discovery window alone would be a minor argument. The
baseline alone would not close the energy chain, because the density term
would remain. The conjunction is the argument.
\end{seamnote}

\section{Where the proposal loses}

A comparison that lists only the wins is not a comparison. Three losses have
to be recorded.

The proposal loses the thermal channel outright, to Sun-Earth L1, and
Chapter 23 shows why: the ground is brighter than the target at every
distance of interest, and a 0.5-metre telescope at 40 kelvin with no regolith
and no night is already funded.

It loses the survey race, because it is not a survey. It does not find
objects; it serves them, at 37 campaigns a year.

And it loses the census, because the arithmetic of Chapter 25 forbids it:
313 years to clear the known 140-metre catalogue. The station is a triage
capability and must be sold as one.

\begin{keytakeaways}
\item The thermal channel is covered twice by funded missions and needs
nothing.
\item The shape channel is covered thinly and the mass channel almost not at
all, and those are the two gaps the proposal addresses.
\item The marginal-value argument puts the alert path first, ahead of any
platform, because it costs almost nothing and closes the documented
routing failure.
\item The proposal loses the thermal channel, loses the survey race and
cannot be a census, and the comparison says so.
\end{keytakeaways}
